\documentclass[11pt,a4paper]{article}

\usepackage{amsmath,amssymb}
\usepackage[utf8]{inputenc}
\usepackage{hyperref}
\usepackage{geometry}
\usepackage{booktabs}
\usepackage{xcolor}
\usepackage{enumitem}
\usepackage{setspace}

\geometry{margin=1in}
\onehalfspacing

\newcommand{\RH}{\ensuremath{\mathrm{RH}}}

\title{%
  What Could Be Built:\\
  Seven Engineering Frontiers Suggested by\\
  the Mathematics of the Cathedral%
}
\author{
  Jason Robert Gochanour
}
\date{April 20, 2026}

\begin{document}

\maketitle

\begin{abstract}
The Cathedral establishes a machine-verified equivalence between
the Riemann Hypothesis and the decay of the Nyman--Beurling
distance, using techniques from spectral theory, harmonic analysis,
variational methods, and number-theoretic sieve constructions.
While the proof itself concerns pure mathematics, the
\emph{structures} it reveals---arithmetic Gram matrices, M\"obius
sieve weights, spectral gaps over multiplicative lattices, and the
Parseval Bridge between time and frequency domains---suggest
concrete engineering possibilities that do not currently exist.
This paper proposes seven grounded, physically realizable
technologies that could be built by adapting the Cathedral's
mathematical insights to engineering domains: (1) prime-spaced
sensor arrays with number-theoretic sidelobe control, (2) arithmetic
compressed sensing with provable recovery guarantees,
(3) M\"obius-weighted adaptive filters for harmonic signals,
(4) machine-verified stability certificates for safety-critical
control systems, (5) prime-gap metamaterials with quasi-random
electromagnetic properties, (6) arithmetic quantum error correction
codes, and (7) formally verified numerical solvers with axiom-tiered
trust. For each proposal, we describe the existing technology, the
mathematical connection to the Cathedral, the engineering design,
and the feasibility assessment.
\end{abstract}

\tableofcontents

% ================================================================
\section{Introduction: From Proof to Possibility}
% ================================================================

The companion paper \emph{From Zeta Zeros to Engineering} mapped
the Cathedral's mathematics to existing engineering concepts---a
dictionary. This paper asks a different question: \emph{what new
things could be built?}

The proposals here are speculative but grounded. Each satisfies
three criteria:

\begin{enumerate}
  \item \textbf{Mathematical foundation}: The proposal is directly
    connected to a mathematical structure in the Cathedral (not
    merely inspired by it).
  \item \textbf{Physical realizability}: The proposed device or
    algorithm could be built with current or near-term technology.
  \item \textbf{Novel capability}: The proposal offers a capability
    that does not exist in current engineering practice.
\end{enumerate}

We are explicit about what is established, what is plausible, and
what is speculative. The Cathedral proves theorems; this paper only
suggests possibilities.

% ================================================================
\section{Prime-Spaced Sensor Arrays}
\label{sec:arrays}
% ================================================================

\subsection{What Exists Today}

Phased array antennas and sonar arrays use either:
\begin{itemize}
  \item \textbf{Uniform spacing} ($\lambda/2$): Simple, but produces
    periodic grating lobes that limit the usable bandwidth.
  \item \textbf{Random spacing}: Eliminates grating lobes but
    produces unpredictable sidelobe levels, requiring more elements
    for reliable performance.
\end{itemize}

Neither fully solves the fundamental tradeoff: predictable pattern
control \emph{and} wide bandwidth.

\subsection{What the Cathedral Suggests}

The Gram matrix $G_N$ encodes the correlation structure of the basis
functions $\{1/(kx)\}$ for $k = 2, \ldots, N$. Its entries depend
on $\gcd(j,k)$ and $\operatorname{lcm}(j,k)$---purely arithmetic
quantities. The key insight: this arithmetic structure is neither
periodic (like uniform spacing) nor random, but \emph{quasi-random
with multiplicative structure}.

The Selberg sieve weights $w_k = \mu(k)(1 - \ln k / \ln N)^+$
provide the optimal beamforming coefficients for this arithmetic
array.

\subsection{The Engineering Design}

\textbf{Proposal}: A sensor array with element positions at
$x_k = k \cdot d$ for $k = 2, 3, 5, 7, 11, \ldots, p_N$ (prime
positions only), with beamforming weights derived from the Selberg
sieve.

\textbf{Properties} (conjectured from the mathematics):
\begin{itemize}
  \item No grating lobes (primes are not periodic).
  \item Predictable sidelobe behavior (controlled by the Gram
    matrix condition number $\kappa \sim N \ln N$).
  \item Optimal weights computable via the Rayleigh--Ritz principle
    applied to the arithmetic Gram matrix.
  \item The Parseval Bridge provides an exact relationship between
    the spatial pattern and the frequency response.
\end{itemize}

\subsection{Feasibility}

\textbf{High.} Non-uniform arrays are already used in radio
astronomy (ALMA, SKA) and military radar. The mathematical
computation (Gram matrix, Selberg weights) is straightforward for
$N \leq 10{,}000$ elements. The novel contribution is using
number-theoretic structure (primes, M\"obius function) instead of
optimization-based aperiodic thinning.

\textbf{Required development}: Simulation of prime-spaced array
patterns; comparison with conventional thinning algorithms;
experimental validation with a small ($N \sim 100$) prototype.

% ================================================================
\section{Arithmetic Compressed Sensing}
\label{sec:cs}
% ================================================================

\subsection{What Exists Today}

Compressed sensing recovers sparse signals from fewer measurements
than the Nyquist rate requires, using random measurement matrices
that satisfy the Restricted Isometry Property (RIP). The standard
approach: use i.i.d.\ Gaussian or Bernoulli matrices, which satisfy
RIP with high probability.

The limitation: random matrices are expensive to store and slow to
apply. Structured matrices (Fourier, Hadamard) are faster but have
weaker RIP guarantees.

\subsection{What the Cathedral Suggests}

The fractional-part functions $\{1/(kx)\}$ form a structured,
deterministic basis on $(0,1)$ with coherence controlled by
$\gcd(j,k)$. The Gram matrix $G_N$ is the ``sensing coherence
matrix,'' and the Cathedral proves its minimum eigenvalue is
strictly positive for all $N$---exactly the condition needed for
stable recovery.

\subsection{The Engineering Design}

\textbf{Proposal}: An arithmetic sensing matrix
$A_{mk} = \{1/(k x_m)\}$ where $x_m$ are measurement points drawn
from $(0,1)$ and $k = 2, \ldots, N$ indexes the ``arithmetic
dictionary.''

\textbf{Recovery guarantee}: If the signal is $s$-sparse in the
arithmetic dictionary, then $O(s \cdot \ln N)$ measurements suffice
for recovery via $\ell_1$ minimization, because the Gram matrix
eigenvalue bound $\lambda_{\min}(G_N) \geq C / \ln N$ provides a
coherence bound.

\textbf{Advantages}:
\begin{itemize}
  \item Deterministic (no randomness needed).
  \item Fast matrix-vector products via FFT-like algorithms
    exploiting the multiplicative structure of the entries.
  \item The sensing matrix has arithmetic structure that natural
    signals (audio, vibration) may already possess: harmonic
    overtones at integer multiples of a fundamental frequency are
    exactly the signals that are sparse in this dictionary.
\end{itemize}

\subsection{Feasibility}

\textbf{Medium.} The mathematical ingredients (Gram matrix
eigenvalue bounds, coherence estimates) exist in the Cathedral.
What's needed: (1)~a formal RIP proof for the arithmetic sensing
matrix, (2)~comparison with Gaussian and Fourier sensing on
synthetic and real signals, (3)~identification of signal classes
that are naturally sparse in the arithmetic dictionary.

\textbf{Most promising application}: Vibration monitoring in
rotating machinery, where the signal is naturally a sum of harmonics
at integer multiples of the rotation frequency.

% ================================================================
\section{M\"obius-Weighted Adaptive Filters}
\label{sec:filters}
% ================================================================

\subsection{What Exists Today}

Adaptive filters (LMS, RLS, Kalman) adjust coefficients to minimize
mean squared error, typically using uniformly weighted or
exponentially weighted observations. These work well for signals
with additive structure (sum of sinusoids, noise).

Signals with \emph{multiplicative} structure---harmonic overtones,
frequency combs, musical signals---are poorly matched by standard
additive models. The harmonic sieve (identifying which overtones
are present) is a combinatorial problem.

\subsection{What the Cathedral Suggests}

The Cathedral's optimal approximation uses M\"obius-weighted
coefficients: $v_k = \mu(k) \cdot w(k)$, where $\mu(k)$ is the
M\"obius function and $w(k)$ is a tapering window. The M\"obius
function naturally encodes multiplicative structure:
\begin{itemize}
  \item $\mu(k) \neq 0$ only for squarefree $k$ (fundamental
    harmonics, no repeated prime factors).
  \item The M\"obius inversion formula $\sum_{d|n} \mu(d) = [n=1]$
    provides exact extraction of fundamental components from a
    harmonic mixture.
\end{itemize}

\subsection{The Engineering Design}

\textbf{Proposal}: An adaptive filter whose weight update rule
incorporates M\"obius structure:
\[
  w_k^{(n+1)} = w_k^{(n)} + \alpha \cdot \mu(k) \cdot e(n) \cdot
  x_k(n),
\]
where $\mu(k)$ pre-weights the update by the M\"obius function,
naturally suppressing composite harmonics and enhancing
fundamental-frequency estimation.

\textbf{Applications}:
\begin{itemize}
  \item \textbf{Musical pitch detection}: Identify the fundamental
    frequency when only overtones are present (the ``missing
    fundamental'' problem).
  \item \textbf{Vibration diagnostics}: Separate bearing defect
    frequencies from shaft harmonics in rotating machinery.
  \item \textbf{Power grid monitoring}: Identify harmonic distortion
    sources in electrical power systems, where harmonics at
    multiples of 50/60 Hz corrupt the signal.
\end{itemize}

\subsection{Feasibility}

\textbf{High.} The M\"obius function is trivially precomputable for
$k \leq 10^6$ (sufficient for audio and vibration). The filter
structure is a minor modification to existing LMS algorithms.
Prototype implementation requires only a few hundred lines of code.

\textbf{Required validation}: Comparison with standard LMS on
synthetic harmonic signals; testing on real audio and vibration
datasets; analysis of convergence properties of the M\"obius LMS.

% ================================================================
\section{Machine-Verified Stability Certificates}
\label{sec:lyapunov}
% ================================================================

\subsection{What Exists Today}

Safety-critical systems (aircraft flight control, nuclear reactor
control, autonomous vehicles) require stability proofs. These are
typically:
\begin{itemize}
  \item Pen-and-paper Lyapunov arguments, checked by human review.
  \item Simulation-based validation (not a proof).
  \item Sum-of-squares (SOS) programming for polynomial systems,
    verified by numerical solvers with floating-point errors.
\end{itemize}

None of these provide the level of assurance that a machine-checked
proof offers.

\subsection{What the Cathedral Suggests}

The Cathedral proves positive definiteness of the Gram matrix
$H_N$ for all $N \geq 1$---a family of stability certificates,
parameterized by $N$, verified by the Lean~4 type-checker. The
proof architecture (axiom registry, tiered assumptions,
\texttt{\#print axioms} discipline) provides a template for
formally verified stability proofs.

\subsection{The Engineering Design}

\textbf{Proposal}: A framework for machine-verified Lyapunov
stability certificates, structured like the Cathedral:

\begin{enumerate}
  \item \textbf{Define} the system dynamics $\dot{x} = f(x)$ in
    Lean~4.
  \item \textbf{Define} a candidate Lyapunov function $V(x)$.
  \item \textbf{Prove} $V(x) > 0$ for $x \neq 0$ (positive
    definiteness).
  \item \textbf{Prove} $\dot{V}(x) < 0$ along trajectories
    (decrease condition).
  \item \textbf{Axiomatize} any numerical estimates (e.g., bounds
    on nonlinear terms) as named, tracked axioms.
  \item \textbf{Audit} the axiom registry to identify which
    stability guarantees depend on which numerical assumptions.
\end{enumerate}

The ``axiom tiering'' system from the Cathedral directly applies:
Tier~1 axioms (kernel) are trusted; Tier~2 axioms (numerical bounds)
are verified by interval arithmetic; Tier~3 axioms (model
assumptions) are validated by testing.

\subsection{Feasibility}

\textbf{Medium-high.} Lean~4 already has the mathematical
infrastructure (Mathlib's analysis library). The gap is
domain-specific: we need formalized control theory definitions
(transfer functions, state-space models, Lyapunov functions) and
connections to numerical verification tools. The Cathedral provides
the organizational template.

\textbf{Impact}: A formally verified autopilot stability proof would
be a major advance in aviation safety certification. Current DO-178C
standards require extensive testing; a machine-checked proof would
complement or replace testing for the most critical properties.

% ================================================================
\section{Prime-Gap Metamaterials}
\label{sec:metamaterials}
% ================================================================

\subsection{What Exists Today}

Metamaterials are engineered structures with electromagnetic
properties not found in nature: negative refractive index, perfect
absorption, cloaking. Their properties depend on the geometry of
sub-wavelength resonant elements arranged in patterns (usually
periodic or graded).

Periodic arrangements produce band gaps (photonic crystals).
Disordered arrangements produce broadband absorption. The tradeoff
between narrow-band performance (periodic) and broadband performance
(disordered) is a fundamental design constraint.

\subsection{What the Cathedral Suggests}

The prime numbers are neither periodic nor random---they are
quasi-random with controlled statistics. The Montgomery--Odlyzko
law shows that zeta zero spacings follow GUE statistics, and the
Cathedral's Gram matrix encodes the correlation structure of
prime-indexed basis functions.

\textbf{Key insight}: A metamaterial with resonant elements at
prime-indexed positions would have:
\begin{itemize}
  \item No Bragg peaks (primes are not periodic)$\to$ no
    narrow-band resonance artifacts.
  \item Controlled density ($\sim 1/\ln N$ primes near $N$)$\to$
    predictable broadband absorption.
  \item GUE-like level repulsion in the resonant spectrum$\to$
    minimum spacing between resonant frequencies, preventing
    destructive interference.
\end{itemize}

\subsection{The Engineering Design}

\textbf{Proposal}: A broadband electromagnetic absorber consisting
of resonant elements (split-ring resonators, patch antennas) with:
\begin{itemize}
  \item Element sizes $a_k$ proportional to $k$ for
    $k = 2, 3, 5, 7, 11, \ldots$  (prime indices only).
  \item Each element resonates at frequency $f_k \propto 1/a_k$,
    producing a ``harmonic comb'' at prime frequencies.
  \item The overall absorption spectrum $A(f)$ is a superposition
    of resonances at prime-indexed frequencies, with the
    Nyman--Beurling theory providing theoretical bounds on the
    spectral coverage.
\end{itemize}

The conjecture: RH implies that this prime-frequency absorber
achieves $A(f) \to 1$ (perfect absorption) as the number of
elements $\to \infty$, because $d_N^2 \to 0$ measures exactly
how well the prime-indexed resonances cover the unit interval.

\subsection{Feasibility}

\textbf{Medium.} Fabrication of metamaterials with non-uniform
element sizes is routine (laser cutting, 3D printing). The novel
contribution is the prime-indexed design rule and the theoretical
connection to the Nyman--Beurling criterion. Simulation (CST
Microwave Studio, COMSOL) should precede fabrication.

\textbf{Key validation}: Simulate the absorption spectrum of a
prime-spaced resonator array and compare with (a)~uniform spacing,
(b)~random spacing, (c)~optimized aperiodic spacing. If the
prime-spaced design offers competitive broadband performance with
a closed-form design rule (no optimization needed), it would be
practically useful.

% ================================================================
\section{Arithmetic Quantum Error Correction}
\label{sec:qec}
% ================================================================

\subsection{What Exists Today}

Quantum error correction (QEC) protects quantum information using
stabilizer codes (surface codes, color codes, toric codes). These
are built from algebraic structures (groups, lattices) over finite
fields $\mathbb{F}_2$.

Current QEC codes do not exploit number-theoretic structure. Their
design is based on topology (surface codes) or algebraic geometry
(Goppa codes), not on multiplicative number theory.

\subsection{What the Cathedral Suggests}

The Gram matrix $G_N$ defines an inner product on the space of
arithmetic functions. Its positive definiteness (proved in the
Cathedral) means that the arithmetic functions $\{1/(kx)\}$ are
linearly independent in $L^2(0,1)$---they form a valid code
basis.

The M\"obius function provides a natural ``error syndrome'':
$\mu(n) = 0$ when $n$ has a squared prime factor (an ``error''),
and $\mu(n) = \pm 1$ when $n$ is squarefree (``no error''). The
M\"obius inversion formula provides exact error decoding.

\subsection{The Engineering Design}

\textbf{Proposal}: A family of quantum codes indexed by $N$, where:
\begin{itemize}
  \item Code qubits correspond to prime indices $p \leq N$.
  \item Stabilizer generators correspond to composite numbers
    $n = p \cdot q$ (products of two primes).
  \item Error syndromes are decoded via M\"obius inversion:
    the error pattern on composites determines the error on primes.
  \item The code distance is bounded by the minimum eigenvalue
    $\lambda_{\min}(G_N) \geq C / \ln N$, giving distance
    $\Omega(\ln N / \ln \ln N)$.
\end{itemize}

\subsection{Feasibility}

\textbf{Low-medium.} This is the most speculative proposal. The
mathematical connection (Gram matrix $\to$ code distance) is
rigorously grounded, but translating from real-valued $L^2$
inner products to binary quantum stabilizers requires a
discretization step that may destroy the number-theoretic structure.

\textbf{Required research}: Formalize the connection between
$L^2(0,1)$ positive definiteness and stabilizer code distance;
construct explicit small codes ($N = 30, 50$) and compute their
parameters; compare with surface codes of similar size.

% ================================================================
\section{Formally Verified Numerical Solvers}
\label{sec:verified}
% ================================================================

\subsection{What Exists Today}

Numerical PDE solvers (finite element, finite difference, spectral
methods) are validated by:
\begin{itemize}
  \item Convergence tests on known analytical solutions.
  \item Mesh refinement studies.
  \item Code-to-code comparisons.
\end{itemize}

None of these constitute a \emph{proof} that the solver produces
correct results for a new problem. The gap between ``tested'' and
``verified'' is the gap between engineering and mathematics.

\subsection{What the Cathedral Suggests}

The Cathedral's architecture---axioms at the boundary, machine-verified
inference in the interior---is directly applicable to numerical
methods:

\begin{enumerate}
  \item \textbf{Axioms as model assumptions}: ``The material is
    linear elastic'' or ``The flow is incompressible'' are stated
    as formal axioms.
  \item \textbf{Discretization as theorem}: ``The FEM solution
    converges to the PDE solution as $h \to 0$'' is proved from
    the axioms using C\'ea's lemma and interpolation estimates.
  \item \textbf{Numerical bounds as tiered axioms}: ``The
    floating-point implementation satisfies $|x_h - x_{\text{exact}}|
    < \varepsilon$'' is a Tier~2 axiom, validated by interval
    arithmetic.
  \item \textbf{Audit trail}: \texttt{\#print axioms solution\_valid}
    shows exactly which assumptions the final result depends on.
\end{enumerate}

\subsection{The Engineering Design}

\textbf{Proposal}: A numerical solver framework in Lean~4 where:
\begin{itemize}
  \item The PDE, boundary conditions, and material model are
    formalized as types.
  \item The discretization (mesh, shape functions, quadrature) is
    formalized with proved convergence.
  \item The linear algebra (assembly, solve) uses verified
    implementations or axiomatized routines with error bounds.
  \item The output is a solution \emph{and} a machine-checked
    certificate that the solution satisfies the stated error bound,
    conditional on the named axioms.
\end{itemize}

\subsection{Feasibility}

\textbf{Medium.} The Cathedral demonstrates that 20,000 lines of
nontrivial mathematics can be machine-verified in Lean~4. A simple
1D FEM solver with formal convergence proof would be $\sim$5,000
lines---well within current capabilities. The challenge is scaling
to 2D/3D problems and connecting to efficient compiled linear
algebra.

\textbf{Near-term milestone}: A formally verified 1D Poisson solver
with proved $O(h^2$) convergence, producing a PDF certificate that
could be submitted as part of an engineering design review.

% ================================================================
\section{Feasibility Summary}
\label{sec:summary}
% ================================================================

\begin{center}
\begin{tabular}{@{}llll@{}}
\toprule
\textbf{\#} & \textbf{Proposal} & \textbf{Feasibility} & \textbf{Timeline} \\
\midrule
1 & Prime-spaced sensor arrays & High & 1--2 years \\
2 & Arithmetic compressed sensing & Medium & 2--4 years \\
3 & M\"obius adaptive filters & High & $<$1 year \\
4 & Verified stability certificates & Medium-high & 2--3 years \\
5 & Prime-gap metamaterials & Medium & 3--5 years \\
6 & Arithmetic quantum error correction & Low-medium & 5--10 years \\
7 & Verified numerical solvers & Medium & 2--4 years \\
\bottomrule
\end{tabular}
\end{center}

The most immediately actionable proposals are \textbf{3} (M\"obius
adaptive filters---a software change to existing LMS algorithms)
and \textbf{1} (prime-spaced arrays---a design rule for existing
hardware). Both could be prototyped in months.

The highest impact proposals are \textbf{4} (verified stability
certificates---transformative for safety-critical systems) and
\textbf{7} (verified solvers---a new standard for engineering
analysis trust). Both require significant Lean~4 library
development but build directly on the Cathedral's organizational
template.

% ================================================================
\section{What This Paper Is Not Claiming}
\label{sec:disclaimer}
% ================================================================

For intellectual honesty:

\begin{enumerate}
  \item We do not claim that the Riemann Hypothesis must be proved
    for any of these technologies to work. The mathematical
    structures (Gram matrices, M\"obius weights, Parseval's theorem)
    exist independently of \RH.
  \item We do not claim that these technologies require the
    Cathedral. An engineer could discover any of these ideas
    independently. The Cathedral's contribution is
    \emph{conceptual}: it demonstrates these structures working
    together in a verified framework.
  \item We do not claim immediate commercial viability. Each
    proposal requires significant research, development, and
    validation before deployment.
  \item We do not claim novelty in the component techniques
    (non-uniform arrays, compressed sensing, adaptive filters).
    The novelty is in the \emph{arithmetic} structure imposed by
    number theory, and in the \emph{verification} framework
    inspired by the Cathedral's proof architecture.
\end{enumerate}

% ================================================================
\section{Conclusion: The Primes as Engineering Material}
% ================================================================

The primes are not a material that engineers can hold. They are not
a force that can be measured. They are not a signal that can be
recorded. But they are a \emph{structure}---a pattern so deep that
it appears in nuclear energy levels, random matrix theory, quantum
chaos, and now in machine-verified proof architectures.

This paper suggests that the prime number structure can be treated
as \emph{engineering material}: a design principle for sensor
arrays, a compression dictionary for harmonic signals, a window
function for spectral analysis, and a template for verified
computation.

The Cathedral did not prove the Riemann Hypothesis. But in the
process of formalizing the attempt, it assembled a toolkit of
mathematical structures that engineers have never used---not
because the structures are unknown, but because the connection
between number theory and engineering has never been made explicit.

These seven proposals are that connection. Some will fail. Some may
succeed. All are grounded in real mathematics, verified by machine.

\begin{quote}
\emph{The primes are not a bridge you can walk across.\\
They are the blueprint for bridges not yet imagined.}
\end{quote}

\vspace{0.5cm}

\begin{thebibliography}{99}

\bibitem{candes2006}
E.~Cand\`es, J.~Romberg, and T.~Tao, \emph{Robust uncertainty
principles: exact signal recovery from highly incomplete frequency
information}, IEEE Trans. Inform. Theory \textbf{52} (2006), 489--509.

\bibitem{haykin2014}
S.~Haykin, \emph{Adaptive Filter Theory}, 5th ed., Pearson, 2014.

\bibitem{khalil2015}
H.~Khalil, \emph{Nonlinear Systems}, 3rd ed., Pearson, 2015.

\bibitem{engheta2006}
N.~Engheta and R.W.~Ziolkowski, \emph{Metamaterials: Physics and
Engineering Explorations}, Wiley, 2006.

\bibitem{gottesman1997}
D.~Gottesman, \emph{Stabilizer Codes and Quantum Error Correction},
PhD thesis, Caltech, 1997.

\bibitem{hesthaven2018}
J.S.~Hesthaven, \emph{Numerical Methods for Conservation Laws},
SIAM, 2018.

\bibitem{cathedral-engineering}
J.R.~Gochanour, \emph{From Zeta Zeros to Engineering: What the
Cathedral's Mathematical Physics Means for Practicing Engineers},
2026.

\end{thebibliography}

\end{document}
